\title{Byte Latent Transformer: Patches Scale Better Than Tokens}

\begin{abstract}
We present the Byte Latent Transformer ({\textsc BLT}), a novel byte-level LLM architecture that, for the first time, achieves parity with tokenization-based LLMs at scale, while also delivering notable gains in inference efficiency and resilience. {\textsc BLT} operates by encoding raw bytes into variable-length patches, which serve as the fundamental computational units. The model adaptively segments these patches according to the entropy of upcoming bytes, thereby dedicating more computation and representational capacity to more complex or unpredictable data regions. We conduct the first scaling analysis of byte-level language models under fixed computational budgets, training models of up to 8B parameters with 4T bytes of input. Our findings affirm that it is viable to scale models using raw byte sequences without reliance on a fixed token vocabulary. Efficiency during training and inference is enhanced because the model can choose larger patches when encountering predictable sequences, further yielding qualitative benefits in reasoning capabilities and generalization across infrequent patterns. On the whole, at constant inference budgets, {\textsc BLT} outperforms tokenization-dependent approaches in scaling effectiveness by increasing both patch and model sizes concurrently.
\end{abstract}

\section{Introduction}
\label{section:intro}

We present the Byte Latent Transformer~(\textbf{BLT}), an architecture that eliminates the need for tokenization by directly processing raw byte sequences. For the first time, BLT achieves comparable results to tokenization-based models when scaled, while demonstrating notable gains in both efficiency and robustness (\S\ref{sec:robustness}). Contemporary large language models (llms) generally train in an end-to-end manner but rely on a tokenization step that heuristically divides byte streams into a predetermined token vocabulary. This approach introduces compression biases that manifest in several limitations, such as heightened sensitivity to specific domains and modalities~\citep{dagangetting}, vulnerability to noisy inputs~(\S\ref{sec:robustness}), limited handling of orthographic variation~\citep{edman2024cute}, and systemic disadvantages for non-dominant languages~\citep{liang2023xlm,petrov2024language,limisiewicz2024myte}.

Tokenization has historically played a crucial role, as training large language models directly on raw bytes has been computationally prohibitive at scale, primarily because of the extended sequence lengths involved~\citep{xue2022byt5}. Some previous studies have addressed this challenge by introducing more efficient forms of self-attention~\citep{el2020characterbert, clark2022canine} or by utilizing architectures that eliminate attention altogether~\citep{wang2024mambabyte}~(\S\ref{section:related_work}). Nonetheless, these advances mainly benefit the training of \textit{smaller models}. In large-scale settings, the main computational bottleneck in Transformers is not the attention component, but rather the sizable feed-forward layers that must be applied to each byte.

To optimize compute allocation, we introduce a dynamic, trainable approach for assembling bytes into \textit{patches}~(\S\ref{section:patching}), along with a novel model architecture that integrates information at both the byte and patch levels. In contrast to tokenization, {\textsc BLT} does not impose a predetermined vocabulary for patches. Instead, flexible byte groupings are transformed into latent patch representations through lightweight, trainable encoder and decoder components. Our experiments demonstrate that this framework leads to \textit{greater} computational efficiency compared to models relying on tokenization.

Traditional tokenization-based llms devote an equal amount of computational resources to each token, prioritizing performance at the expense of efficiency. This approach hinges on compression heuristics for defining tokens, which do not necessarily align with the actual complexity of making predictions for those tokens. Our architecture is fundamentally built on the notion that computational effort should be distributed adaptively, according to demand. For instance, complex computation from a large transformer is generally unnecessary when predicting common suffixes of words, as these predictions usually involve straightforward, low-entropy choices, especially in comparison to initiating a sentence. This principle is operationalized in the design of {\textsc BLT}~(\S\ref{section:architecture}), which incorporates three transformer modules: two compact, byte-oriented \textit{local models} alongside a robust, comprehensive global \textit{latent transformer}~(\autoref{fig:arch_high_level}). 

In order to allocate computational resources dynamically and to decide on how bytes are grouped into patches, {\textsc BLT} leverages the entropy of the next-byte prediction to segment sequences—thus organizing bytes into groups that share similar contextual information density.

We conduct the inaugural flop-constrained scaling analysis of byte-level architectures encompassing up to 8B parameters and 4T training bytes, demonstrating that it is feasible to train large-scale models directly from bytes in an end-to-end manner, foregoing the need for predefined-vocabulary tokenization. In aggregate, {\textsc BLT} achieves parity with Llama 3 regarding training flop-controlled efficiency,\footnote{We estimate a model's computational expenditure by enumerating the number of Floating Point OPerations (flops) required.} while utilizing up to 50\% fewer flops during inference~(\S\ref{section:scaling}). 

Furthermore, our results reveal that operating directly on raw bytes confers meaningful benefits, particularly in capturing the long-tail distribution within the dataset. {\textsc BLT} outperforms conventional tokenizer-based counterparts when encountering noisy input data and displays superior character-level reasoning in areas such as orthographic awareness, phonological patterns, and tasks involving low-resource machine translation~(\S\ref{sec:robustness}).

Additionally, {\textsc BLT} models enable concurrent scaling of both model capacity and patch dimensionality, all while maintaining a constant inference flop budget. Employing longer patch sizes generally leads to computational savings, which can be strategically used to further expand the global latent transformer module, since its execution frequency is reduced. Through a series of inference-flop restricted scaling experiments~(\autoref{fig:fixed_inference_scaling}), we document substantially improved scaling behavior relative to that observed in traditional tokenization-dependent model architectures.

To conclude, this work offers several key advancements: 1) We present {\textsc BLT}, a byte latent LLM framework that adaptively manages computational resources to enhance FLOP efficiency; 2) We demonstrate that our approach achieves training FLOP-parity with Llama 3 up to the 8B parameter level, and enables users to balance slight performance reductions in evaluation metrics against potential FLOP efficiency improvements reaching 50\%; 3) With {\textsc BLT}, large language models gain a novel scaling approach, allowing increases in model size without surpassing a set inference budget; 4) We provide evidence that {\textsc BLT} models show superior resilience to input perturbations and greater sensitivity to sub-word input details frequently overlooked by token-based LLMs. All associated training and inference source code for {\textsc BLT} is made available at~\url{https://github.com/facebookresearch/blt}.

\section{Patching: From Individual Bytes to Groups of Bytes}
\label{section:patching}

Dividing bytes into discrete \textit{patches} enables {\textsc BLT} to flexibly allocate computation depending on contextual information. In Figure~\ref{fig:patching_types}, we present various approaches for segmenting a byte sequence into patches. Formally, a patching function $f_p$ takes an input byte sequence $\pmb{x}=\{x_i,|i=1,\ldots n\}$ of length $n$ and partitions it into $m < n$ patches, forming the sequence $\pmb{p}=\{p_j|j=1,\ldots,m\}$; each $x_i$ is assigned to either 0 or 1, with 1 signifying the initiation of a new patch. For both token-centric and patch-centric models, the predominant driver of computational demand is the number of main Transformer steps required. Within {\textsc BLT}, this corresponds to the quantity of patches into which data is divided for a specific patching function. As such, the patch's average length (referred to as \textit{patch size}) largely dictates data processing cost during both model training and evaluation with a particular patching scheme~(\S\ref{section:flops}).
Subsequently, we introduce three distinct patching strategies: segmenting with a constant byte count per patch~(\S\ref{section:static-patch}), patching according to whitespace boundaries~(\S\ref{section:white-patch}), and adaptive patching guided by entropies from a compact byte language model~(\S\ref{section:dyn-patch}). We conclude by covering incremental patching and contrasting the mechanics of tokenization versus patching~(\S\ref{section:bpe-patch}).

\subsection{Strided Patching Every K Bytes}
\label{section:static-patch}

A common and intuitive approach to segmenting bytes is to divide them into equally sized chunks of length $k$, a method adopted in MegaByte~\citep{yu2023megabyte}. Utilizing a fixed stride simplifies both the implementation of training and inference processes, and provides an uncomplicated way to modify the average patch size, thereby allowing direct control over the computational requirements. Nonetheless, this method has notable drawbacks. Firstly, computational resources are distributed uniformly, without consideration for where they are most required. For example, processing transformer step $j$ may be inefficiently spent if the task is merely to predict whitespace within code, or conversely, insufficient for handling bytes densely packed with significant information, such as mathematical content. Secondly, this technique can yield irregular and context-insensitive segmentation of analogous byte sequences; for instance, identical words may end up fragmented in dissimilar ways.

\subsection{Space Patching}
\label{section:white-patch}

\citet{slagle2024spacebyte} introduces an enhancement to strided patching by generating new patches at each space-like byte\footnote{
    Space-like bytes refer to any byte that does not correspond to a Latin character, digit, or utf-8 continuation byte. Furthermore, every patch must include at least one non space-like byte. }, which frequently align with linguistic unit boundaries in numerous languages. In the Space patching scheme, each word receives an additional latent transformer computation step (i.e., increased computational cost in terms of flops), allowing the model to focus on individual words more accurately. This strategy ensures that the patching of words remains consistent throughout sequences and directs computational resources toward more challenging predictions, which often occur immediately after spaces. For illustration, identifying the first byte in the response to ``Who composed the Magic Flute?\uline{\hspace{.6em}}'' poses a much higher level of uncertainty than subsequent bytes after ``M,'' since the first letter restricts the range of plausible continuations, rendering the rest of ``Mozart'' more predictable. Nonetheless, space patching faces limitations: it is not suitable for all languages and application domains, and crucially, it lacks flexibility in patch size adjustment. To address these issues, we next present a patching approach inspired by the observation that the initial bytes in words tend to be the most challenging to anticipate, while also offering a systematic way to adjust patch sizes.

\subsection{Entropy Patching: Using Next-Byte Entropies from a Small Byte LM}
\label{section:dyn-patch}

Instead of depending on rule-based heuristics like whitespace, we adopt a data-centric strategy to detect instances of high uncertainty in next-byte predictions. We present \textit{entropy patching}, a technique that determines patch boundaries based on entropy calculations.

We train a small byte-level auto-regressive language model on the training data for {\textsc BLT} and compute next byte entropies under the LM distribution $p_{e}$ over the byte vocabulary $\mathcal{V}$:

\begin{equation}
H(x_i) = \sum_{v \in \mathcal{V}} p_{e}(x_i=v|\pmb{x}_{<i}) \log p_{e}(x_i=v|\pmb{x}_{<i})
\end{equation}

We explore two techniques for detecting patch boundaries using entropies $H(x_i)$. The first approach selects points where the entropy exceeds a fixed global threshold, as shown in~\autoref{fig:patching}. The second method detects points that represent a significant increase compared to the preceding entropy value. This latter technique can also be viewed as finding locations where the entropy disrupts an otherwise approximately monotonically decreasing trend within a patch.

\begin{align*}
    \text{Global Constraint}\hspace{1cm}H(x_t)& > \theta_g\\
    \text{Approx. Monotonic Constraint}\hspace{1cm}H(x_t)& - H(x_{t-1}) > \theta_r
\end{align*}

Patch boundaries are determined through a minimal preprocessing procedure carried out while loading the data. This approach contrasts with that of~\citet{nawrot-etal-2023-efficient}, who employ a classifier trained to recognize patch boundaries based on entropy. In our study~(\S\ref{section:experiments}), we evaluate and contrast these two strategies for separating bytes of low and high entropy.

\subsection{The Byte-Pair Encoding (BPE) Tokenizer and Incremental Patching}
\label{section:incremental-patching}
\label{section:bpe-patch}

A wide range of contemporary llms, such as our reference model Llama 3, rely on subword tokenizers like bpe~\citep{gage1994new, sennrich-etal-2016-neural}. In this context, we define ``tokens'' as byte clusters selected from a \textit{finite} vocabulary established before training begins, distinguishing them from ``patches,'' which are composed of dynamically assembled byte sequences lacking a predetermined vocabulary. One key distinction is that, when using tokens, the model cannot directly interact with the fundamental byte-level features.

An important advancement introduced by {\textsc BLT} relative to tokenization-based approaches is its reconceptualization of the balance between vocabulary size and computational requirements. In conventional LLMs, enlarging the vocabulary typically leads to an increase in average token length, resulting in fewer total model steps, but it also necessitates a greater output dimensionality for the model's final projection layer. Consequently, this trade-off constrains tokenization-based methods from realizing substantial changes in token length and the computational expense of inference. For instance, the transition from Llama 2 to Llama 3 raises the average token length from 3.7 bytes to 4.4 bytes, but does so by quadrupling the embedding table’s size.

During generation, {\textsc BLT} must determine at each step whether the current position within the byte sequence coincides with a patch boundary, as this choice dictates whether the Latent Transformer should be activated for additional computation. Critically, this determination must be made based solely on the preceding bytes, without reliance on future sequence elements that have yet to be generated. Therefore, the process of identifying patches must operate without foreknowledge of subsequent bytes, precluding segmentation strategies that depend on future context. In formal terms, a patching scheme $f_p$ is deemed to possess the property of incremental patching if the following holds:
$$f_p(\pmb{x}_{<i}) = f_p(\pmb{x})_{<i}$$
bpe does not qualify as an incremental patching scheme, since identical prefixes may be tokenized in distinct ways contingent on the subsequent sequence, thus violating the property expressed above\footnote{Introducing a dedicated delimiter token at patch boundaries could modify bpe to behave as an incremental patching scheme; however, this results in a longer byte sequence.}.

\section{{\textsc BLT} Architecture}
\label{section:architecture}
{\textsc BLT} consists of a major global autoregressive language model, which processes patch-level representations, in conjunction with a pair of more compact local models responsible for translating byte sequences to patch representations and for reconstructing byte sequences from these patch encodings, respectively~(\autoref{fig:arch_high_level}).

\subsection{Latent Global Transformer Model}

\textit{The Latent Global Transformer} refers to an autoregressive transformer model, denoted as $\mathcal{G}$, comprising $l_{\mathcal{G}}$ layers. This model transforms an input sequence of latent patch representations, $p_j$, into a sequence of output patch representations, $o_j$. Within this work, the index $j$ specifically refers to patches, while $i$ designates individual bytes. The global architecture utilizes a block-causal attention mask~\citep{dubey2024llama}, which enforces that each position can only attend to itself and preceding patches within the same document. Notably, this model accounts for most of the computational overhead (in floating point operations) during both pre-training and inference. Consequently, selecting when to activate this module provides a mechanism for dynamically allocating computational resources across various segments of input and output sequences, dependent on the complexity of those segments.

\subsection{Local Encoder}
\label{section:local}

\textit{The Local Encoder Model}, represented as $\mathcal{E}$, comprises a compact transformer structure utilizing $l_{\mathcal{E}}$ layers, where $l_{\mathcal{E}}$ is substantially fewer than $l_{\mathcal{G}}$. Its principal task is to rapidly translate a series of input bytes, $b_i$, into rich patch-level embeddings, $p_j$. Distinct from standard transformer implementations, each transformer layer in $\mathcal{E}$ is followed by a cross-attention module, which is responsible for aggregating the byte-level information into patch-wise representations~(\autoref{fig:crossattn}).

The pipeline begins by mapping each byte $b_i$ to an embedding vector $x_i$ through a matrix of dimension $\mathbb{R}^{256 \times h_{\mathcal{E}}}$. These embeddings can be supplemented with additional feature information by incorporating hash-embeddings, as detailed in~(\S\ref{sec:hash-emb}). The architecture proceeds with successive applications of transformer and cross-attention components~(\S\ref{sec:enc-cross-attn}), evolving the byte representations into the target patch embeddings $p_i$, which are subsequently consumed by the global transformer module, $\mathcal{G}$.

Within this process, the transformer layers employ a \textit{local block causal} attention mask, permitting each byte to attend to up to $w_{\mathcal{E}}$ previous bytes. This window may cross over dynamically defined patch edges, but never spans document boundaries. The subsequent parts of this section elaborate on the specifics of the embedding mechanisms and the cross-attention structure.

\subsubsection{Encoder Hash n-gram Embeddings}
\label{sec:ngram-emb}
\label{sec:hash-emb}
To build strong and informative representations at every step $i$, it is essential to encode context from the prior sequence of bytes. In {\textsc BLT}, this is realized by representing not only the current byte $b_i$ on its own, but also within the context of a byte n-gram. At every time step $i$, we begin by forming byte-grams

\begin{align}
    g_{i,n}=\{b_{i-n + 1},\ldots, b_i\}
\end{align}

for each byte position $i$ and for $n$ ranging from three to eight.\footnote{
We exclude byte-grams of length $n$ or greater when $i < n$. }

Next, we propose hash $n$-gram embeddings: each possible byte $n$-gram is projected via a hash function into an index within a pre-defined embedding table $E_{n}^{hash}$, corresponding to each $n$ in the set $\{3,4,5,6,7,8\}$~\citep{bai2010learning}. The generated embedding is combined with the original byte's embedding, followed by normalization before serving as input to the local encoder. The enhanced embedding is computed as

\begin{align}
e_i &= x_i + \sum_{n = 3,...,8}E_{n}^{hash}(\text{Hash}(g_{i,n})) \\
\text{where, Hash}(g_{i,n}) &= \text{RollPolyHash}(g_{i,n}) \% |E_{n}^{hash}|
\end{align}

To standardize $e_i$, we divide it by the total count of $n$-gram sizes incremented by one, employing the RollPolyHash technique outlined in Appendix~\ref{appendix:rollpolyhash}. Section~\ref{section:ablations} examines the impact of varying $n$ values and embedding table sizes in $n$-gram hash embeddings on flop-controlled scaling law behavior. Besides hash-based $n$-gram embeddings, we also explored embeddings derived from $n$-gram frequencies; further specifics regarding this investigation can be found in Appendix~\ref{appendix:ngrams}.

\subsubsection{Encoder Multi-Headed Cross-Attention}
\label{sec:enc-cross-attn}
Our approach is largely based on the cross-attention mechanism from the input stage of the Perceiver framework~\citep{jaegle2021perceiver}. A key distinction lies in representing each variable patch by its own latent vector, unlike the fixed latent set used in the original architecture~(\autoref{fig:crossattn}). Additionally, attention computations are restricted solely to the bytes that constitute each respective patch. Specifically, for every patch $p_j$, a corresponding query vector is created by aggregating the byte representations of $p_j$ (using a pooling operation) and then applying a linear transformation, $\mathcal{E}_{C} \in \mathbb{R}^{h_{\mathcal{E}} \times (h_{\mathcal{E}} \times U_{\mathcal{E}})}$, in which $U_{\mathcal{E}}$ denotes the number of encoder cross-attention heads. To express this more precisely, letting $f_{\text{bytes}}(p_j)$ represent the ordered byte sequence associated with patch $p_j$, the computation proceeds as follows:

\begin{align}
P_{0,j} &= \mathcal{E}_{C}(f_\text{bytes}((p_j)), f~ \text{is a pooling function} \\
P_l &= P_{l-1} + W_o\left(\text{softmax}\left(\frac{QK^T}{\sqrt{d_k}}\right)V\right) \\
\text{where } Q_j &= W_q(P_{l-1,j}), K_i = W_k(h_{l-1,i}), V_i = W_v(h_{l-1,i}) \\
h_l &= \text{Encoder-Transformer-Layer}_l(h_{l-1})
\end{align}

Here, $P \in \mathbb{R}^{n_p \times h_{\mathcal{G}}}$ denotes the matrix of $n_p$ patch embeddings that are fed into the global model. These representations are constructed by aggregating the byte embeddings $e_i$ for each corresponding patch $p_j$ through a pooling operation. The matrices $W_q$, $W_k$, $W_v$, and $W_o$ denote the projection layers applied to form the queries, keys, values, and outputs, where keys and values are obtained from projected byte embeddings $h_i$ from the preceding layer (using $e_i$ for the initial layer). To ensure the integrity of patches, a masking approach is employed such that any query $Q_j$ can only attend to those keys and values associated with the bytes from the same patch $j$. Since the model employs multi-headed attention mechanisms over $Q, K$, and $V$, and since the size of patch representations ($h_{\mathcal{G}}$) usually exceeds that of $h_{\mathcal{E}}$, each $P_l$ is represented as several attention heads each of size $h_{\mathcal{E}}$ during cross-attention. These individual head outputs are subsequently concatenated to match the $h_{\mathcal{G}}$ dimensionality. Furthermore, the design incorporates a pre-LayerNorm step before applying the query, key, and value projections, omits positional embeddings within this cross-attention layer, and includes a residual connection enveloping the entire cross-attention block.

\subsection{Local Decoder} 

In parallel with the local encoder, the local decoder $\mathcal{D}$ is implemented as a compact transformer-based architecture, comprising $l_{\mathcal{D}} << l_{\mathcal{G}}$ layers. Its role is to map a sequence of global patch representations $o_j$ back to the original raw bytes $y_i$. Operating autoregressively, the local decoder generates each byte conditioned on the previously decoded outputs. Consequently, it leverages the hidden states produced by the local encoder for the byte sequence as input. The decoding process involves $l_{\mathcal{D}}$ layers in which cross-attention and transformer modules alternate. Specifically, in each layer, cross-attention is performed first—this step derives byte-level representations from the input patch embeddings. Subsequently, the local decoder’s transformer layer processes these byte representations sequentially.

\subsubsection{Decoder Multi-headed Cross-Attention} 

In the decoder cross-attention, the roles of the queries and key/values are interchanged i.e. the byte-representations are now the queries, and the patch representations are now the key/values. The initial byte-representations for the cross-attention are initialized as the byte embeddings from the last encoder layer i.e. $h_{l_{\mathcal{E}}}$. The subsequent byte-representations for layer $l$, $d_{l,i}$ are computed as:

\begin{align}
D_0 &= h_{l_{\mathcal{E}}} \\
B_l &= D_{l-1} + W_o\left(\text{softmax}\left(\frac{QK^T}{\sqrt{d_k}}\right)V\right), \\
  &\text{where } Q_i = W_q(d_{l-1,i}), K_i = W_k(\mathcal{D}_C(o_j)), V_i = W_v(\mathcal{D}_C(o_j)) \\
  D_l &= \text{Decoder-Transformer-layer}_l(B_l)
\end{align}

In this formulation, $W_k$ and $W_v$ refer to the projection matrices for keys and values, respectively, which act through linear transformations in conjunction with the split operation $\mathcal{D}_C$ on the concluding patch representations $o_j$ generated by the global model. The query projection matrix, $W_q$, is applied to the byte-level representations $d_{l-1}$ coming from the preceding decoder transformer layer (or $h_{l_{\mathcal{E}}}$ for the initial layer). The output projection is handled by $W_o$. Accordingly, $B$ has shape $\mathbb{R}^{h_{\mathcal{D}} \times n_b}$, with $n_b$ denoting the total number of output bytes. The subsequent decoder layer representations, $D_l$, are calculated by passing the result of the cross-attention step, $B$, through another decoder transformer layer. Consistent with our approach in the local encoder cross-attention, our implementation employs multi-headed attention, utilizes pre LayerNorms, omits positional embeddings, and incorporates a residual connection surrounding the cross-attention component.

\section{Experimental Setup}
\label{section:experiments}
We meticulously construct a suite of controlled experiments to directly assess the performance of {\textsc BLT} relative to models based on traditional tokenization, with a particular focus on ensuring that {\textsc BLT} does not benefit from potential advantages arising from the use of extended sequence contexts.

\subsection{Pre-training Datasets} The models explored in this study are initially trained on two primary datasets: 1) The Llama 2 dataset~\citep{touvron2023llama}, containing 2 trillion tokens sourced from a diverse array of publicly accessible data, which undergo a rigorous cleaning and filtering procedure to ensure data quality; and 2) {\textsc BLT}-1T, a newly curated collection consisting of 1 trillion tokens accumulated from multiple public sources, including selected subsets of pre-training data from Datacomp-LM~\citep{li2024datacomp}. The Llama 2 dataset forms the basis for scaling law analyses—specifically examining the optimal token count following the guidance of~\cite{dubey2024llama}—and informs the architectural decisions pertaining to {\textsc BLT}. The {\textsc BLT}-1T dataset is employed for an end-to-end pre-training process, facilitating a direct comparison with Llama 3 on subsequent evaluation benchmarks. Both datasets explicitly exclude any content derived from Meta products or associated services. Additionally, in the context of baseline evaluations involving tokenizer-based models, we select the Llama 3 tokenizer (vocabulary size: 128K tokens), which, according to our results, outperforms the Llama 2 tokenizer in establishing robust baseline metrics.

\subsection{Entropy Model} For our experiments, the entropy model is implemented as a byte-level language model, trained on the identical data distribution as the full {\textsc BLT} model. By default, the architecture is a transformer comprising 100M parameters, with 14 layers and a hidden size of 512, as well as sliding window attention that spans 512 bytes. All other hyperparameter choices mirror those used in both our local and global transformers. Variations in model architecture, context window sizes, and parameter counts were examined, as outlined in \autoref{section:ablations}. Notably, when the model’s receptive field is sufficiently limited, the resulting trained entropy model can be represented compactly as a lookup table.

\subsection{Entropy Threshold and Equalizing Context Length} For entropy-based patching approaches, we determine a patching threshold designed to yield a specific average \textit{patch size} over the pretraining data mix. In {\textsc BLT}, as opposed to standard tokenization methods, the \textit{patch size} can be set to any value, substantially affecting the overall context length utilized by the model. To ensure parity in context length and to prevent models with larger patch sizes from receiving undue benefit, we constrain the expected number of bytes per batch to remain fixed. Consequently, when larger patch sizes are used, the sequence length is adjusted downwards. Specifically, for models trained on Llama 2 data, we employ an 8k byte context, whereas for the {\textsc BLT}-1T dataset, we expand the context to an average of 16k bytes, all while keeping the expected batch size at 16M bytes.

Although the mean batch size remains fixed, dynamic patching approaches result in varying proportions of bytes per patch across batches. To optimize computational performance, our {\textsc BLT} training framework composes batches directly from patches, thereby circumventing the need for additional padding within the computationally costly latent transformer. This method guarantees a uniform patch count in each batch. In training, input byte sequences are padded and, when necessary, truncated to 12k bytes for Llama 2 and 24k bytes for {\textsc BLT}-1T datasets. This strategy prevents unexpected memory overhead from particularly large patches within sequences.

\subsection{Entropy Model Context}
\label{section:entropy-context}
Our empirical observations indicate that implementing entropy patching leads to increasingly larger patch sizes within structured formats such as multiple choice evaluations (see \autoref{fig:mmlu_patching} for patching behavior on an MMLU sample), which tend to feature repetitive information. The root of this phenomenon lies in decreased entropy within the repeated content forming the entropy model context. Consequently, during the full-scale application of {\textsc BLT}-Entropy with a patch size of 4.5, we refresh the entropy context via new line insertions and adopt the approximate monotonicity constraint, as this strategy mitigates the “entropy drift” caused by fluctuations in context length. It is important to note that this adjustment pertains solely to our entropy calculations; the method for determining the entropy threshold remains unchanged.

\subsection{FLOPs Estimation} 
\label{section:flops}

Our calculation of transformer floating point operations (flops) closely adheres to the methodology presented in Chinchilla~\citep{hoffmann2022training}, encompassing contributions from feed-forward layers, qkvo projections within the self-attention mechanism, as well as the attention computation and the output projection. One significant distinction in our approach is the treatment of the input embedding layer, which we model as an efficient lookup operation—rather than a dense matrix multiplication—resulting in zero flops attributed to this component. Consistent with prior literature, we approximate the flops for the backward pass as being twice that of the forward pass.

To compute flops \textit{per byte} for {\textsc BLT} models, we add up the flops for the local encoder transformer, the global latent transformer, and the local decoder transformer, together with the cross attention blocks in the encoder and the decoder:

\begin{align}
    \text{FL}_{\text{{\textsc BLT}}} &= \frac{1}{n_p}\,\text{Transf. FL}(h_{\mathcal{G}}, l_{\mathcal{G}}, m = n_{ctx}/n_p, V = 0) \\
   &\quad + \text{Transf. FL}(h_{\mathcal{E}}, l_{\mathcal{E}}, m = w_{\mathcal{E}}, V = 0) \\
   &\quad + \text{Transf. FL}(h_{\mathcal{D}}, l_{\mathcal{D}}, m = w_{\mathcal{D}}, V = 256) \\
    &\quad + \frac{k}{n_p}\,\text{Cross Attn. FL}(h_{\mathcal{E}}, l_{\mathcal{E}}, m = n_p, r = n_p / k) \\
   &\quad + \text{Cross Attn. FL}(h_{\mathcal{D}}, l_{\mathcal{D}}, m = k, r = k/n_p)
\end{align}

In this context, $n_{ctx}$ denotes the sequence length measured in bytes, while $n_p$ indicates the patch size. The parameter $r$ represents the ratio of queries to key/values, and $k$ specifies the ratio of patch dimension to byte dimension; in other words, it signifies how many local model partitions are concatenated to create the overall model representation (for example, $k=2$ in \autoref{fig:crossattn}). The vocabulary size for the output projection, $V$, is utilized exclusively within the local decoder. Depending on whether a module processes a byte-based or patch-based sequence, the attention mechanism adopts a different context length, denoted as $m$. We adjust our attention FLOPs calculations based on the characteristics of each module. Detailed formulas for determining the FLOPs associated with Transformer attention and Cross-Attention are included in Appendix~\ref{section:flops-appendix}.

\subsection{Bits-Per-Byte Estimation} Perplexity only makes sense in the context of a fixed tokenizer as it is a measure of the uncertainty for each token. When comparing byte and token-level models, following previous work~\citep{xue2022byt5, yu2023megabyte, wang2024mambabyte}, we instead report Bits-Per-Byte (BPB), a tokenizer independent version of perplexity. Specifically:

\begin{align}
    \text{BPB}(x)&= \frac{\mathcal{L}_{CE}(\pmb{x})}{\ln(2)\cdot n_{\text{bytes}}}
\end{align}

Here, the total cross-entropy loss reflecting the uncertainty over $\pmb{x}$ is scaled by dividing by both the overall number of bytes in $\pmb{x}$ and a normalization constant.

\subsection{Transformer Architecture Hyperparameters} 

For the transformer modules used in {\textsc BLT}, encompassing both the global and local components, our design primarily mirrors the configuration established in Llama~3~\citep{dubey2024llama}. In the feed-forward networks, the SwiGLU activation~\citep{swiglu} is implemented. For positional encoding within self-attention, we apply rotary positional embeddings (RoPE)~\citep{RoPe}, utilizing a value of $\theta=500000$ as outlined in~\citep{xiong2024effective}. Layer normalization across the model is conducted using RMSNorm~\citep{RMSNorm}. 

Self-attention computations are accelerated with Flash attention~\citep{dao2022flashattention}, specifically when fixed-standard masking schemes such as \textit{block causal} or \textit{fixed-window block causal} are employed; in these cases, a window length of 512 tokens is set for attention masks of a fixed width. Conversely, due to the dynamic nature of patch-dependent cross-attention masking in our model, we leverage Flex Attention\footnote{\url{https://pytorch.org/blog/flexattention}}. This facilitates efficient fused operations, resulting in notably faster training times.

\subsection{{\textsc BLT}-Specific Hyperparameters} 

To evaluate the performance of {\textsc BLT} models, we design experiments focused on two main aspects: scaling behaviors and assessment on downstream tasks. In these studies, model sizes encompass 400M, 1B, 2B, 4B, and 8B parameters. Details of the corresponding architectural hyperparameters for each configuration can be found in Appendix~Table~\ref{tab:arch}.

The initialization of queries for the initial cross-attention layer within the local encoder utilizes max-pooling. Across all {\textsc BLT} model scales, we implement $500,000$ hashes utilizing one hash function, where n-gram lengths are chosen from 3 through 8. A uniform learning rate of $4e-4$ is employed for every model variant. This rate is motivated by a hyperparameter sweep in the range from $1e-3$ to $1e-4$ at 400M and 1B parameter scales, which revealed the optimal value aligns for both token-based and {\textsc BLT} models. For analyses focusing on Llama-2 scaling trends, the training batch sizes used are based on the recommendations of~\cite{dubey2024llama} or an equivalent calculated in bytes.

Regarding optimization, the AdamW optimizer~\citep{adamw} is adopted, with momentum terms $\beta_1 = 0.9$ and $\beta_2 = 0.95$, along with $\epsilon = 10^{-8}$. The learning rate is scheduled with a linear warmup over the first 2000 steps, followed by a cosine decay to zero. Additional regularization and stabilization include a weight decay of 0.1 and global gradient norm clipping capped at 1.0.

\section{Scaling Trends}
\label{section:scaling}

We provide a comprehensive analysis of the scaling behavior observed in byte-level models, offering insights relevant to the continued scaling of {\textsc BLT} architectures. Our investigation into scaling seeks to overcome several shortcomings found in prior studies of byte-level models through the following approaches: (a) conducting comparisons under a compute-optimal training paradigm, (b) training equivalent 8B-parameter models using substantial datasets (reaching up to 1T tokens or 4T bytes) and assessing performance across downstream benchmarks, and (c) examining scaling patterns when inference costs are held constant. Subsequent sections will delve into particular benefits derived from processing byte-sequences.

\subsection{Parameter Matched Compute Optimal Scaling Trends}
\label{section:parameter-matched-scaling}
We conduct experiments on the Llama 2 dataset, training multiple \textit{compute-optimal} bpe and {\textsc BLT} models at four parameter scales, from 1B up to 8B parameters. The analysis involves plotting training flops versus language modeling accuracy on a representative sample of the training mixture. For the bpe models, we adhere to the optimal parameter-to-data ratio established by Llama~3~\citep{dubey2024llama}. This \textit{compute-optimal} approach is theoretically intended to maximize performance for a given training budget on the training set~\citep{hoffmann2022training}, providing a reliable baseline. Each bpe model is paired with a {\textsc BLT} model that is trained on the same data and constructed with a Latent Transformer mirroring the corresponding bpe Transformer's architecture and size.

As shown in~\autoref{fig:scaling-compute-opt} (right), {\textsc BLT} models consistently equal or surpass the performance of their bpe-based counterparts, a pattern that persists across increases in model scale and compute budget. To our knowledge, {\textsc BLT} represents the first Transformer model operating at the byte level to display scaling behavior comparable to BPE architectures when trained under compute-optimal conditions. This observation supports our hypothesis that the optimal balance between model parameters and training compute identified for bpe also extends to {\textsc BLT}, or at minimum, deviates only slightly.

Both enhancements to model architecture and the introduction of dynamic patching play vital roles in achieving bpe scaling behaviors. In ~\autoref{fig:scaling-compute-opt} (left), we present a comparison between models employing space-patching and Llama 3. For SpaceByte \citep{slagle2024spacebyte}, we use the {\textsc BLT} space-patching method while omitting n-gram embeddings and cross-attention, as an approximation. While SpaceByte demonstrates gains over Megabyte, its performance still lags behind that of Llama 3. \autoref{fig:scaling-compute-opt} (right) shows how both architectural innovations and dynamic patching contribute to improved results. When scaled, {\textsc BLT} models achieve comparable performance to leading tokenizer-based models like Llama 3.

Additionally, we note that the selection of tokenizer significantly impacts the results for tokenizer-based models; specifically, models utilizing the Llama-3 tokenizer achieve superior performance compared to those employing the Llama-2 tokenizer when both are trained on identical datasets.

In summary, our {\textsc BLT} model exhibits scaling behaviors that fall between those of Llama 2 and Llama 3 when considerably larger patch sizes are utilized. The average token sizes produced by the bpe tokenizers of Llama 2 and 3 are 3.7 and 4.4 bytes, respectively. By comparison, {\textsc BLT} is able to match these scaling characteristics with average patch sizes of 6 or even 8 bytes. Since inference flop is inversely related to the average patch size, employing an 8-byte patch size results in nearly a 50\% reduction in inference flop. Furthermore, increasing patch sizes appears to enhance model performance as both the dataset and model parameters expand. For example, while {\textsc BLT} with an 8-byte patch size initially underperforms bpe Llama 2 at the 1B scale, it surpasses the bpe version by the time the model reaches the 7B scale. This observation points to the possibility that larger patch sizes may confer even greater advantages at still larger model and data scales, and indicates that the use of even larger patches could become viable as both model capacity and available training resources increase.

\subsection{Beyond Compute Optimal Task Evaluations}
\label{section:evals}
To further investigate scaling behaviors, we train an 8B {\textsc BLT} model surpassing the compute optimal threshold using the larger, higher-quality {\textsc BLT}-1T dataset. We evaluate the model’s capabilities across an array of established classification and generation benchmarks. For this purpose, we focus on a selection of tasks emphasizing common sense reasoning, general knowledge, and code generation abilities:
\paragraph{Classification tasks} encompass ARC-Easy (0-shot)~\citep{Eval_ARC}, Arc-Challenge (0-shot)~\citep{Eval_ARC}, HellaSwag (0-shot)~\citep{Eval_hellaswag}, PIQA (0-shot)~\citep{Eval_piqa}, and MMLU (5-shot)~\citep{hendrycks2020measuring}. Our approach uses prompt-based scoring, estimating the likelihoods assigned to each possible answer, with the reported metric being mean accuracy.
\paragraph{Coding related generation tasks:} We evaluate code generation proficiency using pass@1 on MBPP (3-shot)~\citep{Eval_MBPP} and HumanEval (0-shot)~\citep{Eval_HumanEval}, which test LLMs’ effectiveness at producing correct Python programs.

In \autoref{tab:evals}, we present a comparison among three models that have been trained on the {\textsc BLT}-1T dataset: a bpe Llama 3 tokenizer-based model,\footnote{We select the Llama 3 tokenizer with its 128k token vocabulary, as it outperforms the Llama 2 tokenizer, which has a vocabulary size of 32k.} and two different {\textsc BLT} variants. The first variant applies a space-patching method ({\textsc BLT}-Space), while the second utilizes an entropy-based patching method ({\textsc BLT}-Entropy), both incorporating an approximate monotonicity constraint and resetting the entropy model’s context upon encountering new lines (following the approach described in~\autoref{section:entropy-context}). All models are trained under the same total flop budget. Notably, for the {\textsc BLT}-Entropy model, we further adjust the entropy threshold at inference, lowering it from 0.6 to 0.1, which empirically yields improved task results, albeit at the expense of an increased number of inference steps.

The {\textsc BLT}-Entropy model achieves superior results compared to the Llama 3 model in 4 out of 7 evaluated tasks, even when both models are trained with an equivalent number of bytes. This enhanced performance can be attributed to two primary factors: (1) more efficient utilization of computational resources during training through dynamic patching, and (2) the explicit representation and learning from byte-level data, in contrast to token-level modeling.

Conversely, {\textsc BLT}-Space lags behind the Llama 3 tokenizer across nearly all tasks, except one. Nevertheless, it offers notable savings in inference flops by employing a greater mean patch size of 6 bytes. By contrast, both the bpe and entropy-patching models generate patches averaging about 4.5 bytes within the mixed training dataset. Given identical training resources, the model utilizing the larger patch size is able to process 30\% more data than the other two approaches, a factor which may contribute to {\textsc BLT} deviating further from the compute-optimal regime.

\subsection{Patches Scale Better Than Tokens} With {\textsc BLT} models, it is feasible to scale up both the model capacity and the patch dimensions concurrently without exceeding a fixed computational budget for training and inference or altering the dataset size. Patch-based models uniquely allow patch size to be enlarged arbitrarily, thereby circumventing the efficiency limitations inherent in token-based models with a predetermined vocabulary, as explored in Section~\ref{section:bpe-patch}. Expanding patch size reduces the computational demand, making it possible to allocate these resources toward expanding the global latent transformer, which operates less frequently as a result.

To evaluate whether increasing model size while reducing the number of processing steps on larger input patches can lead to improved performance compared to smaller models with more steps, we carry out a fixed inference scaling study. Our experiments begin with baseline models consisting of 400 million and 3.6 billion parameters, both utilizing the Llama 2 tokenizer. We then identify models that are computation-equivalent with the Llama 3 tokenizer, as well as {\textsc BLT}-Entropy variants that operate on average patch sizes of 6 and 8 bytes within the training datamix (refer to~\autoref{tab:fixed-inf-params} for detailed model configurations). Notably, patch size 8 models are configured with 3 encoder layers, as opposed to the single layer used in smaller patch variants. Each model is trained across a range of prescribed training flop allocations.

As illustrated in \autoref{fig:fixed_inference_scaling}, {\textsc BLT} models demonstrate more favorable scaling behavior compared to tokenization-based counterparts across both inference flop regimes. Notably, BPE models initially exhibit superior performance when training budgets are limited, but they are soon overtaken by {\textsc BLT} models slightly beyond the point of compute-optimality. From a practical standpoint, it may be advantageous to allocate additional resources to pretraining in order to yield a higher-performing model under a constant inference constraint. This is exemplified by 8B-class models—such as Llama 3.1—which have been trained on data volumes approximately two orders of magnitude larger than what would be considered compute-optimal for their parameter count.

The point at which {\textsc BLT} surpasses token-based models has moved somewhat nearer to the compute-optimal threshold as model scale increases, specifically shifting from approximately three times to 2.5 times the compute-optimal budget for the larger flop class models. Likewise, the patch size 8 model in the larger flop regime displays a more rapid scaling curve and surpasses other configurations earlier. As indicated in Section~\ref{section:parameter-matched-scaling}, larger patch sizes tend to yield performance that is more closely aligned with BPE models as the scale of the model increases. We partly ascribe this effect to the declining proportion of the total flops consumed by the byte-level Encoder and Decoder modules, since these components scale less rapidly compared to the Latent Transformer. For example, when expanding overall parameters by a factor of 20—growing from 400M to 8B—{\textsc BLT}'s local model parameters increase by just about twofold. This distinction is significant because increases in patch size influence only the Latent Transformer's flops, not those of the byte-level modules. For this reason, the size ratio of {\textsc BLT}-Entropy ps=8 to Llama 2 rises only modestly from 1.6x to 1.7x as model size increases.

To summarize, our investigation into patch-length scaling reveals that the {\textsc BLT} patch-based model exhibits improved scaling behavior when both patch size and model capacity are expanded concurrently. These positive scaling dynamics continue, and potentially enhance, as model size increases further.

\section{Byte Modeling Improves Robustness} Additionally, we evaluate the robustness of {\textsc BLT} against that of token-based models without intrinsic byte-level data, and introduce a method to adapt existing token-based models to operate on bytes.
\label{sec:robustness}

\subsection{Character-Level Tasks}

An initial incentive for developing byte-level models stemmed from their resilience to noise at the byte level within inputs, as well as their capacity to recognize the underlying elements that form tokens—something that remains a challenge for models relying on tokenization. To assess these attributes, we conduct supplementary experiments utilizing benchmarks designed to test models' robustness against noisy input and their sensitivity to character-level features across English and multiple languages, as well as their handling of digits and phonemes. These findings are detailed in Table~\ref{tab:char_tasks}.

\paragraph{Noisy Data} To assess model robustness, we construct altered versions of the standard classification benchmarks outlined in Section~\ref{section:evals}, introducing various forms of noise. We employ five unique character-level noise manipulations to modify the input text: (a) \textit{AntSpeak}: Transforms text to all uppercase letters, separating each character with a space. (b) \textit{Drop}: Randomly deletes 10\% of the characters within the text. (c) \textit{RandomCase}: Alters the casing of characters such that half are switched to uppercase and half to lowercase, selected at random. (d) \textit{Repeat}: Selects 20\% of characters and repeats them up to four times. (e) \textit{UpperCase}: Converts the entire text to uppercase characters. For evaluation, each noise type is separately applied to the prompt, the completion, or to both, and the scores are averaged across these scenarios. As detailed in Table~\ref{tab:char_tasks}, our experiments on noised versions of HellaSwag~\citep{Eval_hellaswag} reveal that {\textsc BLT} consistently exhibits greater robustness than tokenizer-based baselines, showing an average improvement of 8 points over a model trained on the same dataset and also surpassing the performance of the Llama 3.1 model, which utilizes a substantially larger training corpus.

\paragraph{Phonology - Grapheme-to-Phoneme (G2P)} We evaluate the effectiveness of {\textsc BLT} in converting a sequence of graphemes (alphabetic characters that constitute a word) into its phonemic representation (the pronunciation using phonemes). Table \ref{tab:char_tasks} summarizes the outcomes of the G2P task in a 5-shot scenario on Phonology Bench~\citep{suvarna2024phonologybench}, demonstrating that {\textsc BLT} achieves superior performance compared to the Llama 3 1T tokenizer-based baseline.

\paragraph{CUTE}
To evaluate character-level comprehension, we test {\textsc BLT} using the CUTE benchmark~\citep{edman2024cute}. This suite consists of multiple tasks that fall into three main groups: composition understanding, orthographic similarity analysis, and sequence manipulation skills. The benchmark is notably challenging for most models based on tokenization, which tend to recognize token spellings without being able to leverage this for sophisticated text manipulation. As shown in Table~\ref{tab:char_tasks}, {\textsc BLT}-Entropy surpasses both BPE Llama 3 models on this benchmark by over 25 points. Specifically, our model excels in tasks requiring manipulation at the character level, attaining 99.9\% accuracy for both spelling-focused tasks. These considerable gains are observed even though {\textsc BLT} was trained on just 1/16th the data used for Llama 3.1, underscoring the difficulty BPE-based architectures face in capturing character-level features. Figure \ref{fig:cute_examples} provides examples where the Llama 3 tokenizer-based model underperforms in comparison to our {\textsc BLT} approach. Word deletion and insertion are the rare exceptions where the BPE method does slightly better. While such manipulations may be inherently more challenging for byte-level architectures, the performance difference remains modest, and progressing from characters to words may, in practice, be more tractable than decomposing words into characters. We employ a consistent evaluation protocol for all tasks, utilizing the original Huggingface prompts. Enhanced performance for BPE models might be achieved with further prompt engineering.

\paragraph{Low Resource Machine Translation} We assess the performance of {\textsc BLT} on both directions of translation involving six major language families and twenty-one low-resource languages, each utilizing diverse scripts from the FLORES-101 benchmark~\citep{goyal2022flores}. SentencePiece BLEU scores are presented in Table~\ref{tab:flores}. Our findings indicate that {\textsc BLT} consistently surpasses a model employing the Llama 3 tokenizer, yielding an overall improvement of 2 BLEU points when translating into English and a 0.5-point gain for translations from English. For widely studied language pairs, {\textsc BLT}'s performance is on par with or marginally superior to that of Llama 3. Notably, {\textsc BLT} delivers stronger results across various language pairs within underrepresented language families, highlighting the benefits of byte-level modeling in handling infrequent and diverse byte sequences.

\subsection{Training {\textsc BLT} using Llama 3 Initialization}
\label{sec:distillation}
We investigate an approach where {\textsc BLT} models are able to utilize the parameters of already pre-trained tokenizer-based models in order to enhance training efficiency and speed up convergence. Specifically, we initialize the global transformer parameters of the {\textsc BLT} model with those taken from a pre-trained Llama 3.1 model. After initialization, we proceed by fine-tuning only the global transformer’s weights using a learning rate set to one-tenth of that applied to the local encoder and local decoder, and we continue for the optimal training duration recommended for Llama 3. Results from this procedure are compared to those from a standard {\textsc BLT} baseline in Table~\ref{tab:distill}. The findings indicate that the {\textsc BLT} model initialized from Llama 3.1 considerably surpasses both the standard Llama 3 and the regular {\textsc BLT} baseline under equivalent FLOPs constraints. Further, when contrasted with our {\textsc BLT}-Entropy model reported in Table~\ref{tab:evals}, which was trained on a far larger dataset (1T tokens), {\textsc BLT} from Llama 3.1 still attains better results on the MMLU benchmark. This suggests that our method offers an effective strategy for substantially cutting down on training FLOPs while maintaining high performance.

This approach can alternatively be understood as adapting models that traditionally rely on tokenizers into versions that no longer require them, thereby transforming the pre-trained LLaMA 3.1 architecture into a {\textsc BLT} framework. For a thorough evaluation, we report the results for the original LLaMA 3.1, which was trained on 15T tokens, in Table \ref{tab:distill} and compare these to the {\textsc BLT} model derived from LLaMA 3. In our experiments, the new model shows a modest reduction in performance on the MMLU and HumanEval benchmarks, with a more notable decline across other tasks. These findings indicate that additional efforts are necessary to better utilize the advantages of the pre-trained base model and to enhance performance—particularly through the refinement of training data mixtures and the optimization of hyperparameters.

\section{Ablations and Discussion}
\label{section:ablations}
Here, we analyze the results of ablation studies that support specific architectural decisions in {\textsc BLT}, along with the selected patching method and hyper-parameter configurations utilized for the 8B parameter {\textsc BLT} model trained on the {\textsc BLT}-1T dataset.

\paragraph{Entropy Model Hyper-parameters} In order to examine how adjustments to the size of the entropy model and the length of the context window impact scaling behavior, we develop byte-level entropy transformer models with parameter counts ranging from 1 million up to 100 million, and employ context window lengths spanning from 64 to 512 tokens. The relationships between bpb and training flop scaling laws are visualized in curves generated with our $400m$ and $1b$ {\textsc BLT} models, both trained on the Llama-2 dataset, as depicted in \autoref{fig:entropy_ablation}. Our findings indicate that increases in both the parameter size of the entropy model and the context window length generally enhance scaling performance, though benefits taper off past the 50 million parameter threshold.

\paragraph{Types of Patching}

We conduct an ablation study on the four patching approaches outlined in Section~\ref{section:patching}, specifically: 1) Strided Patching utilizing strides of 4 and 6, 2) Patching at whitespace positions, 3) BPE Tokenizer patching leveraging the Llama 3 tokenizer, and 4) Entropy-driven patching which employs a compact byte-level language model.

Although dynamic patching decreases the actual sequence length, we regulate the sequence lengths to ensure that all patching strategies have a comparable context length. Each model is exposed to an equivalent number of bytes per sequence on average during both training and inference phases, which eliminates the possibility that improved performance is due to handling longer contexts. The outcomes of these control experiments are depicted in Figure~\ref{fig:scaling-compute-opt}. Among the tested methods, all dynamic approaches surpass the static patching baseline, and notably, space patching performs nearly as well as dynamic entropy-based patching.

\autoref{tab:evals_ablation} displays the results of benchmark evaluations for {\textsc BLT} models, contrasting three approaches: models relying on tokenizers, those employing space patching, and those utilizing entropy-based patching. All models were trained on the Llama 2 dataset for an optimized number of training iterations~\citep{dubey2024llama}. While space patching offers a more straightforward approach, avoiding the need to compute entropy dynamically during training, our results show that the improvements associated with entropy-based patching—previously identified in the context of scaling behavior (see Section~\ref{section:scaling})—persist and remain beneficial on downstream benchmark tasks as well.\footnote{The results for space patching reflect earlier experiments that did not incorporate cross-attention; however, the same overall patterns hold when cross-attention is included.}

\paragraph{Cross-Attention} 
In~\autoref{tab:crossattn_ablations_1b}, we systematically remove or alter cross-attention in different locations within the encoder and decoder of {\textsc BLT}. For cross-attention within the encoder, we examine three initialization strategies for the queries: (1) using a uniform learned embedding across all global states, (2) employing a hash embedding derived from the bytes of each patch, and (3) applying a pooling operation over the encoder's hidden states for the bytes in a given patch at the relevant encoder layer.

Our results indicate that integrating cross-attention into the \textit{decoder} yields the highest effectiveness. Employing cross-attention within the encoder offers modest gains, which are apparent only when queries are initialized through pooling. Moreover, cross-attention proves particularly beneficial for the Common-Crawl dataset and demonstrates even greater advantage when utilizing larger patch sizes.

\paragraph{n-gram Hash Embeddings} 

We experiment with n-gram hash embedding vocabularies of sizes 0, 100K, 200K, and 400K, and summarize the findings in Table~\ref{tab:ngram_ablations_1b}. Our observations indicate that hash embeddings yield improvements across every evaluated domain, with Wikipedia and Github benefiting the most (yielding a 0.04 bpb improvement, versus only 0.01 bpb at 15k steps with 8B parameters). Reducing the number of hashes from 500K to 300K at the 8B scale results in a minor performance drop of about 0.001 bpb at 15k steps. This suggests that hash embeddings play a crucial role in achieving parity between {\textsc BLT} and models based on traditional tokenizers; nonetheless, beyond 300K hashes, performance gains taper off. Moreover, our results suggest that the benefits offered by hash embeddings and cross-attention mostly enhance distinct datasets, implying the two methods are largely complementary.

\paragraph{Local Model Hyperparameters} Table \ref{tab:local_layers} presents an ablation study exploring different choices for the number of layers in both the local encoder and decoder. Results indicate that, when using hash n-gram embeddings, {\textsc BLT} achieves strong performance with a very shallow encoder—specifically, a single-layer configuration—while benefiting from a decoder with greater depth.

\section{Related Work}
\label{section:related_work}

\textbf{Character-Level RNNs:} The task of Character Language Modeling has maintained its prominence since the inception of neural architectures~\citep{sutskever2011generating,mikolov2012subword,graves2013generating}, largely due to the capability of these models to handle out-of-vocabulary terms natively, bypassing the need for back-off approaches. In their work, \cite{kim2016character} develop a model that exclusively processes character sequences on the input side, employing convolutional networks alongside highway layers, and passing their outputs to LSTM-based RNNs. This architecture is shown to achieve comparable results to then state-of-the-art RNN language models for English, and it yields superior outcomes on languages with complex morphology, further underscoring the benefits of character-level LLMs. Similarly, \cite{kenter2018byte} adopt byte-level LSTM architectures for machine comprehension tasks and find them to outperform conventional word-level models, notably for morphologically rich languages like Turkish and Russian. In a related direction, \cite{zhang2015character} employ convolutional neural networks built on character-level input for classification problems, and they report improvements over word-level methods in selected contexts. Expanding on these ideas, \citet{chung2019hierarchical} introduce a hierarchical LSTM model incorporating boundary detectors at multiple layers, enabling the system to uncover and exploit latent text hierarchies, leading to enhanced character-level language modeling results. Further, ByteNet~\cite{kalchbrenner2016neural} applies a convolutional neural approach at the character level for machine translation, diverging from the trend of using attention mechanisms.

\textbf{Character-Level Transformers:} Transformer architectures employing attention mechanisms~\citep{vaswani2017attention} in conjunction with subword tokenization strategies~\citep{sennrich-etal-2016-neural} have led to notable advancements in the capabilities of neural models for language modeling and various benchmarks. Nonetheless, the choice of word or sub-word tokens introduces an implicit inductive bias dictating the granularity at which models process language. Seeking to merge the advantages of transformer-based architectures with earlier promising outcomes in character-level modeling, \citet{al2019character} introduced very deep transformers augmented by auxiliary objectives. Their approach yielded transformer-based models that surpassed the performance of prior character-level LSTM language models. Despite these gains, a considerable performance discrepancy remained when compared to word-level large language models. Additionally, GPT-2~\citep{radfordgpt2} reported that, for expansive datasets such as the 1 Billion Word Benchmark, byte-level language models did not reach the effectiveness of word-level counterparts.

\cite{Choe2019BridgingTG} showed that transformer-based LLMs operating at the byte level can surpass subword-level counterparts with a similar number of parameters, yet these byte-level models demand considerably greater computational resources and substantially extended training durations. In a related vein, \cite{el2020characterbert} introduced CharFormer, a BERT-based model that constructs word representations through convolutional operations over character embeddings, yielding improved performance particularly within the medical domain—but again at a significant computational cost. The CANINE model from \cite{clark2022canine} employs 150M parameters and processes raw character sequences without relying on tokens. CANINE’s architecture leverages a deep transformer stack, reminiscent of our global model, combined with a local transformer component and strided convolutions to reduce character-level input length, enabling it to outperform the comparably sized token-level encoder-only model mBERT on multilingual downstream tasks. ByT5~\citep{xue2022byt5} investigated encoder-decoder architectures at the byte level that entirely avoid patching techniques. While their approach demonstrated greater noise robustness and matched the performance of traditional tokenizer-based models using only a quarter of the training data, the absence of patching required attention computations across the full byte sequence, incurring a substantial computational burden. The inherently increased input sequence length associated with byte-level modeling makes efficient scaling particularly challenging. Leveraging the Mamba Architecture~\citep{gu2023mamba}, which enables a constant-sized memory state over very lengthy contexts, \cite{wang2024mambabyte} presented a byte-level implementation of Mamba—again without resorting to patching—that surpasses byte-level transformer models under fixed FLOP budgets at the 350M parameter scale, based on bits-per-byte performance across multiple datasets.

\textbf{Patching-based approaches:} Utilizing patching methods can significantly reduce the otherwise high computational cost associated with byte-level LLMs while maintaining model performance, and a number of studies have reported promising results at limited model and data scales. \cite{nawrot-etal-2022-hierarchical} employ static patching for both downsampling and upsampling, proposing the hourglass transformer architecture, which surpasses competing byte-level models at the 150M parameter scale. Building upon this, \cite{nawrot-etal-2023-efficient} introduce dynamic patching strategies, including an end-to-end learned boundary predictor, a boundary predictor guided by specific tokenizers, as well as an entropy-based patching mechanism comparable to {\textsc BLT}. Their findings indicate that these approaches can surpass standard transformers of the time on language modeling benchmarks at a model size of 40M parameters and over a dataset of roughly 400M tokens. \cite{lester2024training} explore training with sequences compressed via arithmetic coding, which yields higher compression ratios than BPE; using an equal-info windows method, their model exceeds byte-level baselines on language modeling evaluations but does not reach the performance of subword-level baselines.

Our research is primarily influenced by and most closely aligned with MegaByte~\citep{yu2023megabyte}. MegaByte is a decoder-only, causal language model that employs a predetermined, static patching technique for grouping byte sequences and concatenates these representations before inputting them into the model; on the output side, it utilizes a local model within the decoder to reconstruct patches back into bytes. The authors of MegaByte demonstrated that at the 1B parameter level and on a dataset containing approximately 400B bytes, their model performs comparably to models utilizing tokenizers. In our own ablation studies, we systematically compare MegaByte with other approaches and observe that static patching trails compute-optimized, state-of-the-art tokenizer-based models in settings where the number of floating-point operations (FLOPs) is controlled. We further show how {\textsc BLT} succeeds in closing this performance gap. Similarly, \cite{slagle2024spacebyte} reach comparable conclusions regarding MegaByte, and propose augmenting the static patching method by patching not only on bytes but also on whitespaces and analogous characters, in addition to integrating a local encoder module. Their findings indicate enhanced performance relative to token-based transformers in compute-constrained regimes within domains such as Github and arXiv at a 1B parameter scale. We replicate and extend experiments with this model configuration, and our results suggest that further innovations in model architecture are crucial for scaling up byte-level models to reach or surpass the performance of leading token-based approaches like Llama 3.

\section{Limitations and Future Work}

In this study, when making architectural decisions, we train our models using the optimal training steps identified for Llama~3~\citep{dubey2024llama}. It is important to note, however, that these scaling laws were derived specifically for BPE-level transformers and might not yield the ideal ratio of data to parameters for {\textsc BLT}. Establishing scaling laws specifically tailored for {\textsc BLT} is left as future work, as such analysis could potentially reveal even more advantageous scaling properties for our model design. Furthermore, the majority of these experiments were run at a scale of up to 1B parameters. As model size increases to 8B parameters and beyond, the best architectural configurations may shift, potentially enabling further gains at larger model scales.

Current transformer libraries and codebases are primarily optimized for architectures based on tokenizers. Although our work includes theoretically flop-equivalent experiments and leverages efficient implementations like FlexAttention to support layers that differ from standard transformer models, the actual wall-clock performance of our models may still lag behind that of tokenizer-based counterparts. Additional improvements and optimizations could further enhance their efficiency.

Whereas {\textsc BLT} employs an independently trained entropy model specifically for patching, integrating the patching model into a fully end-to-end training framework represents a promising avenue for future investigation. Section \ref{sec:distillation} describes preliminary experiments that suggest the feasibility of ``byte-ifying'' existing tokenizer-based models, such as Llama 3, which have been trained on datasets exceeding 10T tokens. This is achieved by initializing and freezing the global transformer using the pretrained model’s weights. Expanding this line of research could lead to techniques that preserve the advantages of bytefying, while potentially achieving superior performance relative to these tokenizer-based models—without the need for complete retraining from the ground up.

\section{Conclusion}
\label{section:conclusion}

In this work, we introduce the Byte Latent Transformer (\textbf{BLT}), an innovative framework that moves beyond the reliance on static, predefined vocabularies typically required for tokenization in large language models. The central contribution of {\textsc BLT} lies in its adaptive and learnable strategy for forming byte-level patches, which dynamically adjusts computational focus in response to the complexity of the input. This leads to marked gains in efficiency and resilience. Through comprehensive scaling experiments, we show that {\textsc BLT} is able to achieve results comparable to tokenization-based architectures such as Llama 3, particularly at model sizes up to 8B parameters and datasets containing up to 4T bytes. Notably, {\textsc BLT} can achieve up to 50\% savings in inference FLOPs, incurring only minor drops in standard evaluation metrics. Moreover, this approach introduces an expanded avenue for scaling: both model capacity and patch dimensions can be increased together while maintaining a consistent inference cost, making it especially suitable for typical compute constraints in real-world deployments. By processing information directly at the byte level, {\textsc BLT} enhances handling of rare and noisy data instances, leading to better robustness and a more nuanced grasp of fine-grained linguistic elements. Collectively, our findings highlight {\textsc BLT} as a compelling substitute for conventional tokenization pipelines, offering a more adaptable and resource-efficient foundation for future large-scale language models.

\end{document}